\section*{1. Cardinalidad de la consulta}
\addcontentsline{toc}{section}{1. Cardinalidad de la consulta}

Considera las siguientes relaciones:

\vspace{0.5cm}

\begin{center}
\begin{minipage}{0.35\textwidth}
  \centering
  $R =$ \quad
  \begin{tabular}{|c|c|}
    \hline
    \textbf{A} & \textbf{B} \\ \hline
    1 & x \\
    2 & y \\
    2 & z \\
    3 & x \\
    9 & a \\ \hline
  \end{tabular}
\end{minipage}
\begin{minipage}{0.45\textwidth}
  \centering
  $S =$ \quad
  \begin{tabular}{|c|c|c|}
    \hline
    \textbf{B} & \textbf{C} & \textbf{D} \\ \hline
    x & 0 & 3 \\
    y & 2 & 1 \\
    y & 3 & 3 \\
    w & 3 & 0 \\
    y & 4 & 2 \\ \hline
  \end{tabular}
\end{minipage}
\end{center}

\vspace{0.5cm}

Para las siguientes expresiones de \textbf{álgebra relacional}, completa la tabla con el número de tuplas que cada una de ellas produce utilizando las relaciones \textbf{R} y \textbf{S}. Deberás indicar las \textbf{tablas resultantes} en cada caso.

\vspace{0.5cm}

\begin{center}
\begin{tabular}{|l|p{5cm}|}
  \hline
  \textbf{Expresión} & \textbf{Cardinalidad del resultado} \\ \hline
  $R \times S$ & \\ \hline
  $R \bowtie_{D > A} S$ & \\ \hline
  $R \rtimes S$ & \\ \hline
  $R \ltimes S$ & \\ \hline
  $R \bowtie_{A = D} S$ & \\ \hline
  $\rho_{C \leftarrow A}(R) \bowtie S$ & \\ \hline
  $\pi_B(R) - \pi_B(\sigma_{C \ge 3}(S))$ & \\ \hline
  $\pi_A(R) \cap \rho_{A \leftarrow D}(\pi_D(S))$ & \\ \hline
  $\pi_D(S) \bowtie R$ & \\ \hline
  $\gamma_{A; \text{count}(B) \rightarrow t}(R \bowtie S)$ & \\ \hline
\end{tabular}
\end{center}
